\documentclass[11pt]{article}
\usepackage{notes}
\definecolor{AccN}{HTML}{5E1F6A}
\begin{document}
\titlepg{9}{Culmination and Revision:\\One Integrated Framework}{How the eight lectures fit together --- decision framework, annotated formula compendium, misconceptions, the defence-AI pipeline, ethics, and integrative worked problems}
\tableofcontents
\newpage

\section{The Course as One Story}
Lecture~1 established that data can lie and that interpretation errors scale with automation. Lectures 2--3 taught learning from labels (classification, regression, fusion); Lectures 4--5 taught discovery without labels (clustering, anomaly detection); Lectures 6--8 taught decision-making by interaction (single-agent and multi-agent RL). The case studies form a single thread on one notional battlefield: clustering discovered an INFILTRATION archetype (L4); anomaly detection flags such behaviour in real time (L5); a wargame agent learns counter-infiltration positioning (L6); continuous-control swarms provide ISR (L7); and a multi-agent team learns SEAD-first sequencing that denies infiltrators the coverage gaps they exploit (L8).

\section{The Course on One Page}
\begin{center}\begin{tikzpicture}[every node/.style={nd,font=\scriptsize,text width=2.5cm}]
\node[fill=AccN!20,text width=8cm](l1)at(4,3.2){\textbf{L1 Data can lie} --- Simpson, $H_0$, correlation, base rates, missing data, shift};
\node[fill=GrnN!12](l2)at(0,1.6){\textbf{L2 Classify}\\RF, SVM, boosting, metrics};
\node[fill=violet!10](l3)at(2.7,1.6){\textbf{L3 Regress \& fuse}\\OLS, ridge/lasso, Kalman};
\node[fill=teal!12](l4)at(5.4,1.6){\textbf{L4 Cluster}\\$k$-Means, DBSCAN, GMM, PCA};
\node[fill=RedN!10](l5)at(8.1,1.6){\textbf{L5 Detect anomalies}\\IF, OCSVM, autoencoders};
\node[fill=BluN!10](l6)at(1.3,0){\textbf{L6 Q-learning}\\MDP, Bellman, DQN};
\node[fill=GrnN!8](l7)at(4,0){\textbf{L7 Policy gradients}\\PPO, SAC};
\node[fill=AmbN!20](l8)at(6.7,0){\textbf{L8 MARL}\\CTDE, QMIX, MAPPO};
\node[fill=AccN!20,text width=8cm](l9)at(4,-1.5){\textbf{L9--L10 Integrate, validate, govern, present}};
\foreach \x in {l2,l3,l4,l5}{\draw[arr](l1)--(\x);}\draw[arr](l2)--(l6);\draw[arr](l3)--(l7);\draw[arr](l5)--(l8);\draw[arr](l6)--(l7);\draw[arr](l7)--(l8);\foreach \x in {l6,l7,l8}{\draw[arr](\x)--(l9);}
\end{tikzpicture}\\{\small\color{GreyN}Figure 1: concept map of the course.}\end{center}

\section{Master Comparison of Methods}
\begin{center}\scriptsize\renewcommand{\arraystretch}{1.25}
\begin{tabularx}{\textwidth}{>{\bfseries\raggedright\arraybackslash}p{2.5cm}p{2.6cm}p{2.3cm}>{\raggedright\arraybackslash}p{1.9cm}X}\toprule Method & Key assumption & Main knob & Interpret\-able?  & Typical defence use\\\midrule
Logistic regression & linear log-odds & $C=1/\lambda$ & yes (odds ratios) & baseline threat scoring\\
kNN & local similarity, scaled features & $k$ & partly & rare-signature matching\\
Decision tree & axis-aligned splits & depth, $\alpha$ & fully & auditable ROE rules\\
Random forest & decorrelated trees & $B$, $m$ & via importance/SHAP & tabular threat ID\\
Boosting/\allowbreak XGBoost & additive residual fitting & trees, $\eta$, depth & via SHAP & fast embedded classifiers\\
SVM & max margin, kernel & $C$, $\gamma$ & limited & clean high-dim signals\\
Naive Bayes & conditional independence & smoothing & yes & fast baselines, text\\
OLS / ridge / lasso & linear relation & $\lambda$ & yes & planning estimates\\
Kalman filter & linear-Gaussian dynamics & $Q$, $R$ & yes & tracking, fusion\\
$k$-Means & spherical clusters & $k$ & yes (centroids) & behaviour archetypes\\
DBSCAN & density-defined clusters & $\varepsilon$, MinPts & partly & noise / deception\\
GMM & Gaussian components & $K$, covariance & yes & soft membership\\
PCA & linear variance structure & components & loadings & compression, displays\\
Isolation Forest & anomalies isolate easily & $\psi$, trees & partly & real-time screening\\
Autoencoder & normal data compressible & bottleneck & per-feature RE & cyber, sensors\\
DQN family & discrete actions, MDP & $\gamma$, $\varepsilon$, lr & Q-values & wargame advisors\\
PPO / SAC & continuous control & $\varepsilon$, $\alpha$, lr & limited & UAS control\\
QMIX / MAPPO & CTDE, cooperation & as above & limited & team coordination\\\bottomrule
\end{tabularx}\end{center}

\section{Cross-Cutting Ideas}
\subsection{Every method has a complexity knob}
The bias--variance trade-off appears under different names: $k$ in kNN (small $k$, high variance), tree depth and pruning $\alpha$, the number of forest features $m$, SVM's $C$ and $\gamma$, ridge and lasso $\lambda$, the number of boosting trees and learning rate, $k$ and $K$ in clustering, the autoencoder bottleneck, $\gamma$ and $\varepsilon$ in RL, and the PPO clip range. In every case, choose it on validation data, never on the test set.
\subsection{Every method can be fooled by the data}
Survivorship and selection bias corrupt labels (L2--L3) and baselines (L5); scaling choices change clusters (L4) and distances (L2, L5); distribution shift breaks every model, including RL policies moved from simulator to field (L6--L8); confounding creates misleading feature importances and correlations everywhere.
\subsection{Every method needs an honest evaluation}
Per-class metrics and precision at the operational base rate (L2, L5); intervals not points (L3); agreement across methods (L4); alert budgets (L5); many seeds and diverse opponents (L6--L8); confidence intervals and paired tests for comparisons (L1, L10).

\section{Revision Checklist: You Should Be Able To}
\begin{itemize}
\item \textbf{L1:} compute combined rates from counts; run a two-proportion test and a power calculation; compute PPV; explain regression to the mean; recognise survivorship, confounding and shortcut learning.
\item \textbf{L2:} derive the logistic gradient; compute Gini and information gain; explain bagging and boosting; derive the SVM margin; set a cost-optimal threshold; compute precision, recall, $F_\beta$.
\item \textbf{L3:} derive the normal equations; contrast ridge and lasso; perform a Kalman predict--update; build a prediction interval.
\item \textbf{L4:} run $k$-Means by hand; compute a silhouette; label DBSCAN points; do one EM step; compute PCA from a covariance matrix.
\item \textbf{L5:} compute Mahalanobis distance; score with Isolation Forest; set an autoencoder threshold from an FPR budget; design a tiered pipeline.
\item \textbf{L6:} write the Bellman equation; perform a Q-update; explain replay, target networks, Double and Dueling DQN, reward shaping.
\item \textbf{L7:} state the policy-gradient theorem; prove the baseline identity; compute PPO's clipped objective and a GAE advantage.
\item \textbf{L8:} explain CTDE and IGM; compute a COMA advantage; argue autonomy levels under IHL.
\end{itemize}

\section{Practice Examination with Model Answers}
\subsection{Short questions (2 marks each)}
\begin{enumerate}
\item Units A: 30/40 wins in easy terrain, 10/40 in hard; units B: 70/100 easy, 2/20 hard. Overall rates? \emph{A $40/80=50\%$; B $72/120=60\%$ --- yet A wins more in both strata (75 vs 70\%, 25 vs 10\%).}
\item PPV for sensitivity 0.9, specificity 0.95, base rate 1\%? \emph{$0.009/(0.009+0.0495)=15.4\%$.}
\item Gini of a node with class proportions 0.7, 0.3? \emph{$1-0.49-0.09=0.42$.}
\item Scalar Kalman gain for $P=1$, $R=3$? \emph{0.25.}
\item Silhouette with $a=2$, $b=5$? \emph{0.6.}
\item Q-update: $Q=4$, $r=1$, $\max Q'=6$, $\gamma=0.5$, $\alpha=0.5$? \emph{target 4, $\delta=0$, $Q=4$.}
\item PPO term for $r=0.7$, $\hat A=+2$, $\varepsilon=0.2$? \emph{$\min(1.4,\,0.8\times2=1.6)=1.4$, not clipped.}
\item Return of rewards 1, 1, 1 with $\gamma=0.5$? \emph{$1+0.5+0.25=1.75$.}
\item Expected false alerts at FPR 0.2\% on 500,000 normal events? \emph{1,000.}
\item Family-wise error for 10 tests at $\alpha=0.05$? \emph{$1-0.95^{10}=40\%$.}
\end{enumerate}

\subsection{Long question 1 (20 marks)}
\emph{An infiltration-detection system for a border sector must combine radar tracks, seismic sensors and patrol reports. Design it end to end.}

\textbf{Model answer.} Data: tracks with occasional confirmed infiltration labels (survivorship: only detected crossings are labelled), unlabelled seismic signals, free-text reports. Methods: Kalman-based fusion of radar and seismic detections into tracks; clustering of track behaviour to maintain archetypes (L4), refreshed monthly; Isolation Forest plus autoencoder on track features for novel behaviour (L5), thresholds set from an alert budget; a supervised classifier where labels exist, with class weights and per-class recall; an RL or MARL patrol-allocation advisor trained in simulation (L6--L8). Evaluation: precision at the true base rate, detection delay, confidence intervals, red-team crossings. Governance: Level-2 recommendations, analyst confirmation, audit logs; privacy safeguards for civilian movement.

\subsection{Long question 2 (20 marks)}
\emph{Critically evaluate the claim: ``Our MARL agents beat expert human planners, so we should deploy them to command drone swarms.''}

\textbf{Model answer.} Statistical: 82.4\% vs 81.3\% on one evaluation is within sampling error (a z-test on 300 vs 150 missions gives $p\approx0.8$). External validity: simulation results typically lose 15--20 points in the field without domain randomisation; human planners adapt across scenarios while the agents were trained on one. Adversarial: a fixed policy is exploitable; learned messages can be spoofed; emergent degenerate behaviours (holding fire) reveal reward mis-specification. Legal and ethical: IHL distinction and proportionality require accountable human judgement; penalties are not safeguards. Conclusion: use the agents as Level-2 planning advisors with human approval, after staged validation, red-teaming and audit design.

\section{The Decision Framework}
\begin{enumerate}
\item \textbf{Check the data first} (Lecture~1): how was it collected, what is missing, what is the base rate, which confounders drove assignment, does the deployment domain match?
\item \textbf{Do I have labels?} Yes $\to$ supervised: discrete output $\to$ classification (random forest strong default for tabular data; decision tree for auditability; SVM for clean, moderate-size, high-dimensional data; boosting for speed and rare-class recall); continuous output $\to$ regression (boosting for non-linear tabular data; ridge, lasso or elastic net for interpretable models; SVR for robustness; Kalman-based fusion for tracking). No $\to$ unsupervised: find groups ($k$-Means first, DBSCAN for shape and noise, GMM for soft membership, Ward for a visual hierarchy) or flag the unusual (Isolation Forest to screen, autoencoder to confirm, sequence or graph models for behaviour and networks).
\item \textbf{Is the problem sequential and interactive?} Yes $\to$ RL: small discrete actions $\to$ Double/Dueling DQN; continuous actions $\to$ PPO (stable default) or SAC (robust, sample-efficient).
\item \textbf{Several agents?} Yes $\to$ MARL with CTDE: MAPPO default, QMIX for cooperative discrete tasks, MADDPG for continuous or mixed settings.
\item \textbf{Validate and govern}: per-class metrics, confidence intervals, deployment-domain tests, explanation, autonomy level, audit trail.
\end{enumerate}

\section{Annotated Formula Compendium}
\subsection{Data and statistics (L1)}
\begin{itemize}
\item $\text{Error}=\text{Bias}^2+\text{Variance}+\sigma^2$ --- complexity trades bias for variance; noise is irreducible.
\item $P(\ge1\text{ false positive})=1-(1-\alpha)^m$; Bonferroni $\alpha/m$ --- many tests manufacture discoveries.
\item $\text{PPV}=\dfrac{sp}{sp+(1-c)(1-p)}$ --- alarms are mostly false when the base rate $p$ is small.
\item $n\approx(z_{1-\alpha/2}+z_{1-\beta})^2[p_1q_1+p_2q_2]/\Delta^2$ --- compute sample size before a trial.
\item $\mathbb E[X_2|X_1]=\mu+\rho(X_1-\mu)$ --- regression to the mean.
\item $D_{KS}=\max_x|F_1(x)-F_2(x)|$ --- detect distribution shift.
\end{itemize}
\subsection{Supervised learning (L2--L3)}
\begin{itemize}
\item $P=\frac{TP}{TP+FP}$, $R=\frac{TP}{TP+FN}$, $F_\beta=(1+\beta^2)\frac{PR}{\beta^2P+R}$ --- choose $\beta$ by error costs.
\item $G=1-\sum p_k^2$; $H=-\sum p_k\log_2p_k$ --- tree split criteria.
\item $\text{Var}_{\text{ens}}=\rho\sigma^2+(1-\rho)\sigma^2/B$ --- forests need decorrelated trees.
\item SVM: margin $2/\|w\|$; $\min\frac12\|w\|^2+C\sum\xi_i$; $K(x,x')=e^{-\gamma\|x-x'\|^2}$.
\item Threshold $p>C_{FP}/(C_{FP}+C_{FN})$ --- cost-optimal decision with calibrated probabilities.
\item $\hat\beta=(X^\top X)^{-1}X^\top y$; ridge $(X^\top X+\lambda I)^{-1}X^\top y$; lasso soft-threshold $\text{sign}(z)\max(|z|-\lambda/2,0)$.
\item Kalman: $K=P^-H^\top(HP^-H^\top+R)^{-1}$; $\hat x=\hat x^-+K(z-H\hat x^-)$; $P=(I-KH)P^-$; scalar $K=P/(P+R)$.
\end{itemize}
\subsection{Unsupervised learning (L4--L5)}
\begin{itemize}
\item $J=\sum_j\sum_{x\in C_j}\|x-\mu_j\|^2$; $s(i)=\frac{b-a}{\max(a,b)}$.
\item $r_{ik}=\frac{\pi_k\mathcal N_{ik}}{\sum_j\pi_j\mathcal N_{ij}}$; BIC $=p\log N-2\log L$.
\item PCA: $Cv=\lambda v$, EVR $=\lambda_j/\sum\lambda$.
\item $D_M=\sqrt{(x-\mu)^\top\Sigma^{-1}(x-\mu)}$; $s(x,\psi)=2^{-E[h(x)]/c(\psi)}$; $\text{RE}=\|x-\hat x\|^2$.
\end{itemize}
\subsection{Reinforcement learning (L6--L8)}
\begin{itemize}
\item $G_t=\sum\gamma^kr_{t+k+1}$; horizon $\approx1/(1-\gamma)$.
\item $Q^*=\sum P[R+\gamma\max Q^*]$; $Q\leftarrow Q+\alpha[r+\gamma\max Q'-Q]$.
\item Double DQN: $y=r+\gamma Q(s',\arg\max Q(s';\theta);\theta^-)$; dueling $Q=V+A-\bar A$; shaping $F=\gamma\Phi(s')-\Phi(s)$.
\item $\nabla J=\mathbb E[\nabla\log\pi\,\hat A]$; $\hat A=r+\gamma V(s')-V(s)$; $L^{\text{CLIP}}=\mathbb E[\min(r_t\hat A,\text{clip}(r_t)\hat A)]$; SAC adds $\alpha H(\pi)$.
\item QMIX $\partial Q_{\text{tot}}/\partial Q_i\ge0$; MAPPO $\frac1N\sum L_i^{\text{CLIP}}$; COMA $A_i=Q-\sum\pi_i(a')Q(s,(a',\mathbf a_{-i}))$.
\end{itemize}

\section{Misconceptions, Explained}
\begin{enumerate}
\item \textbf{``More data always helps.''} More of the same biased data reinforces the bias; Afghan mine data could not fix Iraq performance.
\item \textbf{``High accuracy means a good model.''} Under imbalance, predicting the majority class scores highly and detects nothing; use per-class recall and precision at the operational base rate.
\item \textbf{``Percentages can be added or averaged.''} Totals come from counts (Simpson).
\item \textbf{``Not significant means no effect.''} Under-powered tests miss real effects; report confidence intervals.
\item \textbf{``DBSCAN needs $k$.''} It discovers the number of clusters from density.
\item \textbf{``A baseline biases the policy gradient.''} It has zero expected contribution and only reduces variance.
\item \textbf{``Replay is only for efficiency.''} It also breaks correlation; without it DQN often diverges.
\item \textbf{``t-SNE distances are meaningful.''} Only local neighbourhoods are preserved.
\item \textbf{``QMIX just sums Q-values.''} VDN sums; QMIX mixes non-linearly and monotonically, conditioned on the state.
\item \textbf{``Q-values are win probabilities.''} They are expected discounted rewards.
\item \textbf{``Simulation success equals field success.''} Expect 15--20-point gaps without domain randomisation and human-opponent testing.
\end{enumerate}

\section{The Integrated Defence-AI Pipeline}
\textbf{Stage 1 --- ingestion and data quality:} radar, IR, ELINT, SIGINT and OSINT feeds; bias audits, shift detection, label-noise estimates, missing-data strategy. \textbf{Stage 2 --- threat identification:} classification, regression and Kalman-based fusion. \textbf{Stage 3 --- pattern discovery:} clustering of behaviours, dimensionality reduction for display, tiered anomaly detection. \textbf{Stage 4 --- decision support:} DQN advisors for discrete courses of action, PPO/SAC for continuous control, MARL for teams. \textbf{Stage 5 --- human--AI interface:} explanations, calibrated confidence, Level-2 approval, ROE and IHL checks, audit logs, override. The pipeline is iterative: anomalies update baselines, outcomes retrain classifiers, operator feedback updates every model. Typical latencies: edge inference 10--100 ms; decision support 0.1--1 s; human decision 1--10 s; sensor-to-action 1.5--11 s.

\section{Ethics: Eight Principles}
(1) Meaningful human control --- Level 2 for lethal decisions. (2) Data integrity --- audits and monitoring. (3) Explainability --- every recommendation needs a rationale. (4) Proportionality --- a legal judgement no reward can encode. (5) Distinction --- not reliably solvable by AI at operational tempo. (6) Adversarial robustness --- red-team before deployment. (7) Accountability --- a full chain from input to outcome. (8) Honesty about simulation --- overstating capability is itself an ethical failure.

\section{Integrative Worked Problems}
\begin{example}{a cyber-defence pipeline for 10 million packets a day}
\textbf{Data:} audit labelled incidents for survivorship bias (only detected attacks are labelled); stratify by network segment. \textbf{Screening:} an Isolation Forest at 0.002 ms/packet costs 20 s/day; at FPR 1.1\% it flags about 110,000 normal packets. \textbf{Confirmation:} an autoencoder on flagged traffic at FPR 0.6\% leaves about 660 false alerts, still too many; aggregate to per-host alerts and add a supervised classifier for known attack families. \textbf{Base rate:} with 0.1\% attacks, report precision at that rate, not ROC-AUC. \textbf{Human role:} analysts triage ranked alerts; containment actions such as isolating a host require approval. \textbf{Validation:} replay held-out weeks, red-team novel attacks, monitor drift.
\end{example}
\begin{example}{interpreting a surprising exercise result}
Units using a new AI planning aid won 62\% of engagements against 80\% for others. Before concluding anything: tabulate counts by terrain and threat level (Simpson); check whether AI units were assigned harder sectors; compute standardised rates; check sample sizes and confidence intervals; look for regression to the mean if the aid was given to previously weak units. In the Lecture~1 example the aid was better in every terrain.
\end{example}
\begin{example}{from simulation to deployment for a drone swarm}
SAC achieved 96\% in simulation. Estimate the transfer gap from analogous results (PPO 94\% $\to$ 78\% without randomisation); apply domain randomisation (wind, noise, mass, latency); validate in hardware-in-the-loop and range trials with confidence intervals; wrap the policy in a certified safety envelope (geofence, separation, fallback); deploy at Level 2 for any armed role.
\end{example}

\clearpage
\part*{Part C --- Deeper Integration}
\addcontentsline{toc}{part}{Part C --- Deeper Integration}

\section{How the Methods Connect}
The course has introduced some forty methods. They are far fewer ideas than that. This section shows the mathematical threads that tie them together; understanding these links is the best preparation for an unseen problem.

\subsection{One neuron, three names}
A logistic-regression classifier computes $\sigma(w^\top x+b)$ (Lecture~2). A single sigmoid neuron computes exactly the same thing (Lecture~1, deep learning). A softmax policy over discrete actions computes $\text{softmax}(Wx)$, the multi-class version, and its REINFORCE gradient $(\mathbf 1_a-\pi)\,x\,G$ has the same form as the cross-entropy gradient $(\mathbf 1_y-p)\,x$ (Lecture~7). Supervised classification and policy-gradient RL differ only in what multiplies the gradient: a correct label in one case, a return in the other. That is why a policy can be pre-trained by behavioural cloning and then fine-tuned by RL with no change of architecture.

\subsection{$k$-means is a special case of EM}
In a Gaussian mixture with equal spherical covariances $\sigma^2I$ and equal weights, the E-step responsibilities are $r_{ik}\propto\exp(-\|x_i-\mu_k\|^2/2\sigma^2)$. As $\sigma\to0$ the largest term dominates and every responsibility becomes 0 or 1 --- assign each point to its nearest centroid. The M-step then sets each mean to the average of its points. That is Lloyd's algorithm. GMMs are ``soft $k$-means''; $k$-means is ``hard EM''.

\subsection{Inverse-variance weighting is everywhere}
Fusing two independent estimates $\hat x_1$ (variance $P_1$) and $\hat x_2$ (variance $P_2$) gives $\hat x=(\hat x_1/P_1+\hat x_2/P_2)/(1/P_1+1/P_2)$ with variance $(1/P_1+1/P_2)^{-1}$. This single formula is:
\begin{itemize}
\item the scalar \textbf{Kalman update} (Lecture~3), with gain $K=P_1/(P_1+P_2)$;
\item \textbf{Bayes' rule} for two Gaussians, which is why the Kalman filter is an optimal Bayesian filter;
\item the \textbf{belief update} of a POMDP (Lecture~6) in the linear-Gaussian case;
\item the fixed-effect \textbf{meta-analysis} formula that combines studies in medicine and psychology.
\end{itemize}
Wherever two uncertain sources must be combined, ask what their variances are.

\subsection{Regularisation is a prior}
Maximising the posterior $p(w\mid D)\propto p(D\mid w)p(w)$ with Gaussian noise and a Gaussian prior $w\sim\mathcal N(0,\tau^2I)$ gives
\[-\ln p(w\mid D)=\frac1{2\sigma^2}\sum_i(y_i-w^\top x_i)^2+\frac1{2\tau^2}\|w\|^2+\text{const},\]
which is \textbf{ridge regression} with $\lambda=\sigma^2/\tau^2$. A Laplace prior $p(w)\propto e^{-|w|/b}$ gives \textbf{lasso}. Strong regularisation means a confident prior that coefficients are small; weight decay in neural networks and the KL penalty in PPO and SAC are the same idea: stay close to a reference unless the data insist.

\subsection{From Bellman to DQN in five steps}
(1) The Bellman optimality equation defines $Q^*$. (2) With a model, apply it as an update: \textbf{value iteration}. (3) Without a model, replace the expectation by one sampled transition: \textbf{Q-learning}. (4) Replace the table by a function with parameters: \textbf{semi-gradient Q-learning}. (5) Stabilise with replay and a frozen target: \textbf{DQN}. Every step is an approximation of the one before; each new failure mode (sample noise, generalisation error, the deadly triad) is the price of that approximation.

\subsection{Boosting is gradient descent on functions}
Gradient descent updates parameters, $\theta\leftarrow\theta-\eta\nabla_\theta\mathcal L$. Gradient boosting updates the \emph{prediction function}: it computes the negative gradient of the loss at each training point, $-\partial\ell/\partial F(x_i)$ (the residual, for squared error), fits a tree to those values, and adds a small multiple of the tree, $F\leftarrow F+\eta T$. The learning rate, early stopping and regularisation of boosting are the learning rate, early stopping and regularisation of gradient descent.

\subsection{Attention is kernel regression}
Nadaraya--Watson kernel regression predicts $\hat y(x)=\sum_iK(x,x_i)y_i/\sum_jK(x,x_j)$. Self-attention computes $\sum_i\text{softmax}_i(q^\top k_i/\sqrt d)\,v_i$ --- a kernel-weighted average of values, with the kernel $\exp(q^\top k/\sqrt d)$ learned through the query and key projections. A Transformer is a stack of learned, data-dependent smoothers.

\subsection{PCA is a linear autoencoder}
An autoencoder with one hidden layer of $q$ linear units, trained to minimise squared reconstruction error, learns to project onto the same subspace as the first $q$ principal components. Non-linear autoencoders generalise PCA; their reconstruction error is the anomaly score of Lecture~5, just as PCA reconstruction error is.
\begin{center}\begin{tikzpicture}[font=\scriptsize]
\tikzset{cmb/.style={nd,minimum width=2.5cm,minimum height=0.6cm}}
\node[cmb,fill=BluN!8](lr)at(0,0){logistic regression};\node[cmb,fill=BluN!8](nn)at(4.2,0){neural network};\node[cmb,fill=RedN!8](pg)at(8.4,0){policy network};
\node[cmb,fill=GrnN!10](km)at(0,-1.3){$k$-means};\node[cmb,fill=GrnN!10](gmm)at(4.2,-1.3){GMM / EM};\node[cmb,fill=AmbN!15](kf)at(8.4,-1.3){Kalman / Bayes};\node[cmb,fill=AmbN!15](po)at(12.2,-1.3){POMDP belief};
\node[cmb,fill=GrnN!10](pca)at(0,-2.6){PCA};\node[cmb,fill=BluN!8](ae)at(4.2,-2.6){autoencoder};\node[cmb,fill=RedN!8](an)at(8.4,-2.6){anomaly score};
\node[cmb,fill=RedN!8](be)at(0,-3.9){Bellman equation};\node[cmb,fill=RedN!8](ql)at(4.2,-3.9){Q-learning};\node[cmb,fill=RedN!8](dq)at(8.4,-3.9){DQN};
\draw[arr](lr)--node[above,font=\tiny]{stack}(nn);\draw[arr](nn)--node[above,font=\tiny]{returns}(pg);
\draw[arr](gmm)--node[above]{$\sigma\to0$}(km);\draw[arr](gmm)--node[above,font=\tiny]{Bayes}(kf);\draw[arr](kf)--node[above,font=\tiny]{hidden}(po);
\draw[arr](pca)--node[above]{non-linear}(ae);\draw[arr](ae)--node[above]{error}(an);
\draw[arr](be)--node[above]{sample}(ql);\draw[arr](ql)--node[above]{network}(dq);
\end{tikzpicture}\\*{\small\color{GreyN}Figure C1: some of the threads that connect the course's methods.}\end{center}

\section{One Air-Defence Sector, Worked End to End}
This case uses the tools of every lecture on a single, numerical problem. Work through it before looking at each answer.
\subsection{The situation}
A sector processes 2,000 radar tracks a day, of which about 4 are hostile. A classifier (Lecture~2) flags hostile tracks with recall 0.95 at a false-positive rate of 1\%. An independent ESM sensor detects hostile emitters with probability 0.8 and false-alarms on 10\% of other tracks.
\subsection{Step 1 --- What does an alarm mean? (Lecture~1)}
Classifier alarms per day: $4\times0.95=3.8$ true and $1{,}996\times0.01=19.96$ false; precision $=3.8/23.8=0.16$. Five out of six alarms are false. Requiring ESM confirmation as well: true $3.8\times0.8=3.04$, false $19.96\times0.1=2.00$; precision $0.60$, at the cost of missing 20\% of the threats the classifier caught. A two-stage design trades recall for precision; the commander must decide whether 0.76 hostiles a day lost to the second stage is acceptable, or whether unconfirmed alarms go to a human for a quick look instead of being discarded.
\subsection{Step 2 --- Where is it? (Lecture~3)}
For a confirmed track, radar gives range 52.0 km (variance 4 km$^2$) and an IR-derived estimate gives 49.0 km (variance 16). Fused range $=(52/4+49/16)/(1/4+1/16)=51.4$ km, variance $3.2$ km$^2$ (standard deviation 1.8 km): closer to the radar, more certain than either.
\subsection{Step 3 --- When should the system cue an engagement? (Lectures~1--2)}
If a missed hostile costs 50 times a false cue, the cost-optimal threshold on the classifier's calibrated probability is $p^*=C_{FP}/(C_{FP}+C_{FN})=1/51\approx0.02$: cue on anything above 2\%. The threshold follows from the costs, which are a command judgement --- and from calibration, without which ``2\%'' means nothing.
\subsection{Step 4 --- Is something new happening? (Lectures~4--5)}
Over a week the daily mean radar cross-section of unidentified tracks drifts upward; a CUSUM on that statistic alarms on day~5. Clustering the week's unidentified tracks with HDBSCAN finds a small dense group with an unfamiliar combination of speed and altitude --- a candidate new platform type. It is passed to analysts, not to the engagement chain.
\subsection{Step 5 --- What should be done? (Lectures~6--8)}
A wargame-trained advisory agent proposes three courses of action with estimated values and uncertainty from an ensemble of critics. The commander sees the options, the assumptions and the agent's track record against diverse opponents, and decides. Expected-utility reasoning (Lecture~1) makes the value judgements explicit; the agent never selects lethal action on its own (Lecture~8).
\subsection{Step 6 --- How do we know it works? (Lectures~9--10)}
Evaluate each component on later data from the deployment environment; report recall of the hostile class with a confidence interval (with only a few hostiles a day, months of data are needed for a tight one --- Lecture~10); red-team the classifier with novel and deceptive tracks; monitor drift continuously; document limitations and the conditions under which human-only procedures resume.

\section{Practice Examination II}
Each question carries the marks shown. Model answers follow the questions.
\begin{enumerate}
\item \textbf{[6]} A new targeting aid was used in 200 engagements and succeeded in 114; without it, 142 of 200 succeeded. By difficulty: easy engagements 81/90 with the aid versus 136/170 without; hard engagements 33/110 with the aid versus 6/30 without. Is the aid better or worse? Explain with counts.
\item \textbf{[4]} A scalar Kalman filter has prior variance 9 and measurement variance 3. Give the gain and the posterior variance.
\item \textbf{[4]} A point has mean intra-cluster distance 1.2 and mean nearest-other-cluster distance 3.0. Compute its silhouette and interpret.
\item \textbf{[4]} An Isolation Forest with $\psi=256$ gives a point a mean path length of 5. Compute its score ($c=10.24$) and interpret.
\item \textbf{[4]} Q-learning: $Q(s,a)=2$, $r=1$, $\max_{a'}Q(s',a')=5$, $\gamma=0.95$, $\alpha=0.2$. Update $Q(s,a)$.
\item \textbf{[4]} PPO with $\varepsilon=0.2$: ratio $0.7$, advantage $-3$. Compute the clipped objective term and say whether the gradient pushes the ratio further down.
\item \textbf{[6]} A report claims ``our MARL team beat the human team, 84\% versus 80\%, over 100 engagements each''. Assess the claim.
\item \textbf{[8]} Propose, with justification, a method for each: (a) grouping 50,000 unlabelled intercepts; (b) predicting fuel use per sortie with intervals; (c) a drone that must hold position in gusts; (d) flagging unusual logins.
\end{enumerate}
\subsection*{Model answers}
\begin{enumerate}
\item Overall the aid looks worse: $114/200=57\%$ against $142/200=71\%$. Within each stratum it is better: easy $81/90=90\%$ versus $136/170=80\%$; hard $33/110=30\%$ versus $6/30=20\%$. The reversal (Simpson's paradox) arises because the aid was used mostly in hard engagements (110 of its 200) while the comparison group fought mostly easy ones (170 of 200). The aid is better; the totals mislead because they mix different case loads. Compare within strata, or standardise both groups to the same mix of difficulty.
\item $K=9/(9+3)=0.75$; posterior variance $=9\times3/12=2.25$.
\item $s=(3.0-1.2)/3.0=0.60$: reasonably well placed in its cluster.
\item $s=2^{-5/10.24}=0.71$: shorter paths than average; likely anomalous, worth review.
\item $Q=2+0.2(1+0.95\times5-2)=2+0.2(3.75)=2.75$.
\item Unclipped $0.7(-3)=-2.1$; clipped $0.8(-3)=-2.4$; PPO takes the minimum, $-2.4$, which is the clipped (constant) term, so there is no gradient pushing the ratio further below $1-\varepsilon$ --- the update stops reducing this action's probability.
\item Difference 4 points; standard error $\approx\sqrt{0.84(0.16)/100+0.80(0.20)/100}=0.054$; $z\approx0.74$, $p\approx0.46$: not demonstrated. Also ask about opponent types, simulation fidelity, and whether the humans had the same information and interface.
\item (a) Embed or engineer features, then HDBSCAN or GMM, validated by stability and analyst review. (b) Gradient boosting or regularised regression with quantile or conformal prediction intervals; evaluate on later sorties. (c) SAC or PPO with domain randomisation, wrapped in runtime assurance; or a classical controller with learned gain scheduling. (d) Per-user and peer-group baselines with an Isolation Forest or autoencoder, thresholds set by the analysts' alert budget.
\end{enumerate}

\section{Notation Used in the Course}
\begin{center}\small\renewcommand{\arraystretch}{1.25}
\begin{tabularx}{\textwidth}{>{\bfseries}l X}\toprule Symbol & Meaning\\\midrule
$x,\;y,\;\hat y$ & input vector, true target, prediction\\
$w,\;b,\;\theta$ & weights, bias, all parameters\\
$\mathcal L,\;\ell$ & total loss, per-example loss\\
$\alpha$ (learning) / $\alpha$ (tests) & learning rate / significance level --- context decides\\
$\lambda$ & regularisation strength; also eligibility-trace decay in TD($\lambda$)\\
$\sigma(\cdot)$ & sigmoid function; $\sigma$ alone is a standard deviation\\
$\mu,\;\Sigma$ & mean vector, covariance matrix\\
$P,\;Q,\;R,\;K,\;F,\;H$ (filtering) & state covariance, process noise, measurement noise, Kalman gain, transition and measurement matrices\\
$s,\;a,\;r,\;s'$ & state, action, reward, next state\\
$\gamma,\;G_t$ & discount factor, return\\
$\pi,\;V,\;Q,\;A$ (RL) & policy, state value, action value, advantage\\
$\varepsilon$ & exploration rate; PPO clip range; noise term --- context decides\\
$r_t(\theta)$ & PPO probability ratio $\pi_\theta/\pi_{\text{old}}$\\
$Q_{tot}$ & team value in value-factorisation MARL\\\bottomrule
\end{tabularx}\end{center}
The same letter means different things in different fields ($Q$ is an action value in RL and a process-noise covariance in Kalman filtering). In an examination, define every symbol the first time you use it.

\section{Reading a Research or Evaluation Paper Critically}
\begin{enumerate}
\item What exactly is the task, and is it the operational task or a convenient proxy?
\item Where did the data come from, and how might it differ from deployment data?
\item What is the baseline, and is it a fair one (tuned with the same effort)?
\item Are results averaged over several seeds or splits, with uncertainty shown?
\item Is the metric appropriate for the class balance and the cost of errors?
\item Could information have leaked from test to training (random splits of sequential data, preprocessing on all data)?
\item Were many variants tried and only the best reported (the garden of forking paths)?
\item What happens outside the tested conditions --- does the paper say?
\item Would the method's failure be graceful or catastrophic, and is a human in the loop?
\end{enumerate}

\section{Examination Technique}
Use numbers rather than adjectives (``RF 97.5\%, bomber recall 96.1\%, 2.8 ms''); state limitations (sample size, confidence intervals, simulation, base rate); connect lectures (the L4 $\to$ L5 $\to$ L6 $\to$ L8 thread); address the autonomy level and accountability in every deployment answer; show worked arithmetic (Kalman gain, silhouette, Q-update, PPV); and be suspicious of absolutes such as ``always'', ``proves'' or ``guarantees''.
\end{document}
